% ============================================================
%  PROPOSAL LOMBA - WANI (WhatsApp AI Native Integration)
%  Platform Omnichannel Berbasis AI untuk UMKM Indonesia
% ============================================================
\documentclass[12pt,a4paper,oneside]{article}

% ── Encoding & Font ──────────────────────────────────────────
\usepackage[utf8]{inputenc}
\usepackage[T1]{fontenc}
\usepackage{times}

% ── Bahasa ───────────────────────────────────────────────────
\usepackage[bahasa]{babel}
\usepackage{indentfirst}
\setlength{\parindent}{1.5cm}

% ── Geometry ─────────────────────────────────────────────────
\usepackage[top=3cm, bottom=3cm, left=3.5cm, right=3cm]{geometry}

% ── Packages ─────────────────────────────────────────────────
\usepackage{graphicx}
\usepackage{xcolor}
\usepackage{hyperref}
\usepackage{enumitem}
\usepackage{listings}
\usepackage{booktabs}
\usepackage{tabularx}
\usepackage{caption}
\usepackage{setspace}
\usepackage{ragged2e}
\usepackage{float}
\usepackage{csquotes}
\usepackage{fancyhdr}
\usepackage{titlesec}
\usepackage{multicol}
\usepackage{makecell}
\usepackage{textcomp}

% ── Colors ───────────────────────────────────────────────────
\definecolor{primary}{HTML}{0F766E}     % warm teal
\definecolor{accent}{HTML}{D97706}      % amber
\definecolor{dark}{HTML}{1E293B}
\definecolor{lightbg}{HTML}{F8FAFC}
\definecolor{teallight}{HTML}{CCFBF1}

% ── Hyperlinks ───────────────────────────────────────────────
\hypersetup{
    colorlinks=true,
    linkcolor=primary,
    urlcolor=primary!70!black,
    citecolor=dark
}

% ── Page style ───────────────────────────────────────────────
\setlength{\headheight}{14pt}
\pagestyle{fancy}
\fancyhf{}
\fancyhead[L]{\small\textcolor{primary}{\textbf{WANI}} -- Proposal Lomba}
\fancyhead[R]{\small\textcolor{dark}{\thepage}}
\renewcommand{\headrulewidth}{0.4pt}
\renewcommand{\headrule}{\hbox to\headwidth{\color{primary}\leaders\hrule height 0.4pt\hfill}}

% ── Section formatting ───────────────────────────────────────
\titleformat{\section}
  {\Large\bfseries\color{primary}}{}{0em}{}[\vspace{-0.5ex}\rule{\textwidth}{0.5pt}]
\titleformat{\subsection}
  {\large\bfseries\color{dark}}{}{0em}{}
\titleformat{\subsubsection}
  {\normalsize\bfseries\color{dark!70}}{}{0em}{}

\titlespacing{\section}{0pt}{2.5ex}{1.5ex}
\titlespacing{\subsection}{0pt}{1.5ex}{0.8ex}

% ── Spacing ──────────────────────────────────────────────────
\onehalfspacing

% ── Listing style ────────────────────────────────────────────
\lstdefinestyle{bash}{
    language=bash,
    basicstyle=\footnotesize\ttfamily,
    backgroundcolor=\color{lightbg},
    frame=single,
    rulecolor=\color{primary!30},
    breaklines=true,
    postbreak=\mbox{\textcolor{primary}{$\hookrightarrow$}\space},
    aboveskip=1em,
    belowskip=1em
}

\lstdefinestyle{output}{
    basicstyle=\footnotesize\ttfamily,
    backgroundcolor=\color{lightbg},
    frame=single,
    rulecolor=\color{gray!30},
    breaklines=true,
    aboveskip=1em,
    belowskip=1em
}

% ── Boxed content ────────────────────────────────────────────
\newcommand{\highlightbox}[2][primary!10]{%
    \noindent\colorbox{#1}{%
        \parbox{\dimexpr\linewidth-2\fboxsep}{\textbf{#2}}%
    }%
    \vspace{0.5em}%
}

% ── Two-column list ──────────────────────────────────────────
\newenvironment{twocol}{%
    \begin{multicols}{2}[]
    \begin{itemize}[leftmargin=1em, nosep]
}{%
    \end{itemize}
    \end{multicols}
}

% ── References ───────────────────────────────────────────────
\usepackage{cite}

\begin{document}

% ============================================================
%  COVER PAGE
% ============================================================
\begin{titlepage}
    \centering

    % Logo placeholder
    % \includegraphics[width=3.5cm]{logo-unida.png}\\[1.5cm]

    \vspace*{1cm}

    {\color{primary}\rule{0.7\textwidth}{1.5pt}}\\[1.5em]
    {\LARGE\bfseries\color{dark}
        WANI\\[0.3em]
        {\large\textit{WhatsApp AI Native Integration}}\\[1.2em]
    }
    {\Large\bfseries\color{primary}
        Platform Omnichannel Berbasis AI\\[0.2em]
        untuk Digitalisasi UMKM Indonesia\\[1.5em]
    }
    {\color{primary}\rule{0.7\textwidth}{1.5pt}}\\[2em]

    {\large\bfseries PROPOSAL}\\ [0.5em]
    {\large Diajukan untuk Mengikuti \textbf{[Nama Lomba]}}\\[3em]

    \begin{center}
        \begin{tabular}{lc}
            \textbf{Ketua Tim} & : Farrel Ghozy Affifudin \quad ([NIM]) \\[0.6ex]
            \textbf{Anggota 1} & : [Nama Anggota 1] \quad ([NIM]) \\[0.6ex]
            \textbf{Anggota 2} & : [Nama Anggota 2] \quad ([NIM]) \\[0.6ex]
            \textbf{Anggota 3} & : [Nama Anggota 3] \quad ([NIM]) \\[0.6ex]
        \end{tabular}
    \end{center}

    \vfill

    \begin{center}
        \textbf{PROGRAM STUDI TEKNIK INFORMATIKA}\\[0.3ex]
        \textbf{FAKULTAS [Nama Fakultas]}\\[0.3ex]
        \textbf{UNIVERSITAS DARUSSALAM GONTOR}\\[0.3ex]
        \textbf{2026}
    \end{center}
\end{titlepage}

% ============================================================
%  RINGKASAN EKSEKUTIF
% ============================================================
\newpage
\section*{Ringkasan Eksekutif}
\addcontentsline{toc}{section}{Ringkasan Eksekutif}

\noindent
\textbf{WANI} (\textit{\textbf{W}h\textbf{A}tsapp \textbf{N}ative \textbf{I}ntegration} / \textit{WhatsApp Niaga}) adalah platform omnichannel terintegrasi yang menghubungkan \textit{WhatsApp bot} bertenaga kecerdasan buatan, \textit{REST API server}, \textit{admin dashboard}, dan \textit{static website generator} dalam satu kesatuan sistem yang kohesif. Platform ini dirancang khusus untuk menjawab tantangan digitalisasi Usaha Mikro, Kecil, dan Menengah (UMKM) di Indonesia.

\vspace{0.8em}
\noindent
WANI menawarkan empat solusi inti:

\begin{enumerate}[nosep]
    \item \textbf{AI Chatbot WhatsApp} -- Otomatisasi layanan pelanggan 24/7 dengan \textit{intent detection} berbasis \textit{large language model} (LLM) yang mampu memahami dan merespons 6 kategori intensi pelanggan (order, inquiry, greeting, complaint, unknown, escalate).
    \item \textbf{Multi-Layer Security Pipeline} -- Sistem pertahanan injeksi 3-tier (regex $\rightarrow$ ML classifier $\rightarrow$ deep LLM judge) yang melindungi dari \textit{prompt injection}, kebocoran PII, dan eksfiltrasi data -- menjadikannya salah satu sistem chatbot AI paling aman untuk lingkungan UMKM.
    \item \textbf{Dashboard Administrasi Terpusat} -- Manajemen produk, pesanan, pelanggan, dan konfigurasi sistem dalam satu antarmuka web modern (React 19, TypeScript 6, Tailwind CSS 4).
    \item \textbf{Website Generator Otomatis} -- Pembuatan situs web statis UMKM dari data toko dan produk, lengkap dengan integrasi WhatsApp (\textit{click-to-order}), tanpa perlu keahlian teknis.
\end{enumerate}

\vspace{0.5em}
\noindent
Dibangun di atas \textit{stack} teknologi modern dan stabil (Bun, Express 5, Prisma 7, PostgreSQL, React 19, Astro 6.4) dengan metodologi Extreme Programming, WANI telah melalui tahap pengembangan yang matang dan siap untuk diimplementasikan sebagai solusi digitalisasi bagi pelaku UMKM Indonesia.

\vfill

% ============================================================
%  TABLE OF CONTENTS
% ============================================================
\newpage
\tableofcontents
\newpage

% ============================================================
%  1. PENDAHULUAN
% ============================================================
\section{Latar Belakang}

\subsection{Darurat Digitalisasi UMKM Indonesia}

Usaha Mikro, Kecil, dan Menengah (UMKM) merupakan pilar fundamental perekonomian Indonesia. Data Kementerian Koperasi dan UKM mencatat bahwa sektor UMKM menyumbang \textbf{61\%} terhadap Produk Domestik Bruto (PDB) nasional (setara Rp9.580 triliun) dan menyerap sekitar \textbf{97\%} dari total tenaga kerja \cite{kemkop2023}. Namun, kesenjangan digital masih menjadi hambatan serius.

Berdasarkan survei Kementerian Koperasi dan UKM tahun 2022 \cite{kemkop2022survei}, hanya sekitar \textbf{20\%} UMKM yang telah mengadopsi teknologi digital secara aktif dalam operasional bisnis mereka. Sebagian besar masih mengelola komunikasi pelanggan, pencatatan pesanan, dan data produk secara manual -- menggunakan buku catatan, spreadsheet sederhana, atau bahkan hanya mengandalkan ingatan. Kondisi ini menyebabkan:

\begin{itemize}
    \item \textbf{Inefisiensi operasional} -- Waktu respons lambat, pesanan terlewat, dan kesalahan pencatatan.
    \item \textbf{Hilangnya peluang bisnis} -- Pelanggan beralih ke kompetitor yang lebih responsif.
    \item \textbf{Ketidakmampuan analisis} -- Tidak ada data historis untuk pengambilan keputusan bisnis.
    \item \textbf{Ketiadaan \textit{web presence}} -- UMKM kesulitan membangun kehadiran \textit{online} karena keterbatasan biaya dan keahlian teknis.
\end{itemize}

\subsection{WhatsApp: Kanal Komunikasi Utama, Namun Terabaikan}

WhatsApp telah menjadi infrastruktur komunikasi \textit{de facto} bagi UMKM Indonesia. Dengan \textbf{90,9\%} tingkat penetrasi sebagai platform media sosial paling banyak digunakan di Indonesia \cite{wearesocial2024}, WhatsApp digunakan secara luas untuk menerima pesanan, menjawab pertanyaan, dan membangun hubungan dengan pelanggan. Ironisnya, potensi besar ini belum dimanfaatkan secara optimal:

\begin{itemize}
    \item Percakapan bisnis bercampur dengan chat pribadi.
    \item Tidak ada sistem \textit{ticketing} atau pelacakan pesanan.
    \item Tidak ada integrasi dengan manajemen inventaris dan produk.
    \item Bot WhatsApp yang ada di pasaran umumnya tidak memiliki keamanan AI yang memadai.
\end{itemize}

\subsection{Ancaman Prompt Injection pada Chatbot AI}

Seiring dengan maraknya adopsi \textit{large language model} (LLM) untuk layanan pelanggan, muncul ancaman serius berupa \textit{prompt injection attack}. OWASP \cite{owasp2025} mengidentifikasikan \textit{prompt injection} (LLM01:2025) sebagai kerentanan nomor satu pada aplikasi LLM. Pelaku UMKM yang menggunakan chatbot AI tanpa lapisan keamanan yang memadai berisiko mengalami:

\begin{itemize}
    \item Kebocoran data sensitif (data pelanggan, informasi bisnis).
    \item Manipulasi sistem untuk menjalankan perintah berbahaya.
    \item Eksfiltrasi \textit{system prompt} dan konfigurasi internal.
\end{itemize}

\subsection{Solusi: WANI}

Dari permasalahan di atas, lahirlah \textbf{WANI} (\textit{\textbf{W}h\textbf{A}tsapp \textbf{N}ative \textbf{I}ntegration} / \textit{WhatsApp Niaga}) -- sebuah platform omnichannel terintegrasi yang menghubungkan \textit{WhatsApp Bot} bertenaga AI, \textit{REST API server}, \textit{dashboard} administrasi, dan \textit{website generator} statis. WANI dirancang khusus untuk membantu UMKM Indonesia mengelola komunikasi pelanggan, manajemen produk dan pesanan, serta membangun \textit{web presence} secara otomatis -- dengan keamanan AI berlapis yang menjadi standar tertinggi di kelasnya.

% ============================================================
%  2. KEUNGGULAN DAN INOVASI
% ============================================================
\section{Keunggulan dan Inovasi}

WANI menawarkan sejumlah inovasi yang membedakannya dari solusi \textit{chatbot} dan \textit{point-of-sale} yang sudah ada di pasaran:

\begin{table}[H]
\centering
\caption{Perbandingan dengan Solusi Lain}
\label{tab:perbandingan}
\begin{tabularx}{\textwidth}{p{3cm}Xp{3.5cm}p{2.5cm}}
\toprule
\textbf{Aspek} & \textbf{WANI} & \textit{Chatbot} \textbf{Konvensional} & \textbf{SaaS POS} \\
\midrule
AI Security & 3-tier guardrail (regex $\rightarrow$ ML $\rightarrow$ Judge) & Tidak ada / basic & Tidak relevan \\
\textit{Intent Detection} & 6 kategori (Zod validated) & Rule-based terbatas & Tidak ada \\
\textit{Omnichannel} & WhatsApp + Dashboard + Web & Chat only & POS only \\
\textit{Source Code} & \textbf{Open Source} (\href{https://github.com/FarrelGhozy/WANI}{GitHub}) & Proprietary & Proprietary \\
\textit{Data Ownership} & \textbf{100\% milik pengguna} (self-hosted) & Tersimpan di pihak ketiga & Tersimpan di pihak ketiga \\
\textit{Web Presence} & Auto-generate website UMKM & Tidak ada & Tidak ada \\
\textit{Biaya} & Gratis (self-hosted) + API LLM saja & Berlangganan bulanan & Berlangganan bulanan \\
\textit{Latest Stack} & Bun, Express 5, React 19, Prisma 7 & Bervariasi, sering usang & Proprietary \\
\bottomrule
\end{tabularx}
\end{table}

\subsection{Inovasi Utama}

\begin{enumerate}
    \item \textbf{18-Langkah AI Pipeline dengan Keamanan Berlapis} \\
    Tidak seperti \textit{chatbot} konvensional yang hanya meneruskan pesan ke LLM, WANI memproses setiap pesan melalui 18 langkah terstruktur yang mencakup normalisasi input, deduplikasi, \textit{rate limiting}, pemindaian PII, pertahanan injeksi 3-tier, \textit{output grounding}, dan redaksi PII. Pipeline ini memastikan setiap respons AI aman, akurat, dan kontekstual.

    \item \textbf{Sistem Pertahanan Injeksi 3-Tier Pertama di Kelasnya} \\
    WANI mengimplementasikan tiga lapis pertahanan yang bekerja secara berjenjang:
    \begin{itemize}
        \item \textbf{Tier 1 (wajib, $\sim$0ms):} Deteksi berbasis 9 kelompok pola regex, termasuk deteksi homoglif Cyrillic/Greek dan normaliasi leetspeak.
        \item \textbf{Tier 2 (kondisional, $\sim$500ms):} Klasifikasi ML melalui OpenRouter model cepat untuk kasus yang mencurigakan.
        \item \textbf{Tier 3 (kondisional, $\sim$1500ms):} Analisis mendalam dengan konteks percakapan penuh untuk kasus batas (\textit{borderline}).
    \end{itemize}

    \item \textbf{Architecture Modular 4-in-1} \\
    WANI terdiri dari empat sub-proyek independen yang dapat dikembangkan, diuji, dan di-\textit{deploy} secara terpisah: \texttt{api/} (Express 5 REST server), \texttt{dashboard/} (React 19 SPA), \texttt{wa-bot/} (Baileys 6 WhatsApp client), dan \texttt{web-gen/} (Astro 6.4 site generator). Setiap sub-proyek menggunakan TypeScript dengan \texttt{verbatimModuleSyntax} dan \texttt{erasableSyntaxOnly}.

    \item \textbf{Open Source dan Data Ownership} \\
    Seluruh kode sumber WANI tersedia secara \textit{open source} di \href{https://github.com/FarrelGhozy/WANI}{GitHub}. Data bisnis sepenuhnya berada di kendali pengguna (\textit{self-hosted} PostgreSQL), tanpa ketergantungan pada pihak ketiga.

    \item \textbf{Website Generator untuk Web Presence} \\
    WANI tidak hanya mengelola komunikasi dan transaksi, tetapi juga secara otomatis menghasilkan situs web statis UMKM dari data toko dan produk -- lengkap dengan tombol ``Pesan via WhatsApp'' yang terintegrasi.
\end{enumerate}

% ============================================================
%  3. ARSITEKTUR DAN CARA KERJA
% ============================================================
\section{Arsitektur dan Cara Kerja}

\subsection{Arsitektur Sistem}

WANI mengadopsi arsitektur modular tiga-komponen yang saling terhubung melalui protokol HTTP dengan autentikasi Bearer token:

\begin{figure}[H]
\centering
\begin{lstlisting}[style=output,language={}]
+------------------+      HTTP/JSON       +------------------+
|    Dashboard     | <------------------> |   API Server     |
|    React 19      |     port 5173        |   Express 5      |
|    Vite 8        |   (Vite Proxy)       |   port 3001      |
+------------------+                      +--------+---------+
                                                    |
                                    +---------------+---------------+
                                    |  Bearer Auth  |  Bearer Auth  |
                                    |  POST /api/qr | POST /api/chat|
                                    |  GET /api/qr  | PUT /api/store|
                                    +---------------+---------------+
                                                    |
                                          +---------v----------+
                                          |     WA Bot         |
                                          |   Baileys 6        |
                                          |   Prisma 7         |
                                          |   PostgreSQL       |
                                          +--------------------+

          +------------------+
          |    Web-Gen       |
          |   Astro 6.4      |
          |   Static Site    |
          +------------------+
\end{lstlisting}
\caption{Arsitektur Sistem WANI}
\label{fig:arsitektur}
\end{figure}

\subsection{Alur Data Utama}

\begin{enumerate}
    \item \textbf{Pelanggan mengirim pesan} via WhatsApp $\rightarrow$ WA Bot menerima pesan.
    \item WA Bot meneruskan pesan ke \texttt{POST /api/chat} pada API Server.
    \item API Server menjalankan \textbf{AI Pipeline 18 langkah}: normalisasi $\rightarrow$ rate limit $\rightarrow$ PII scan $\rightarrow$ 3-tier injection defense $\rightarrow$ LLM completion $\rightarrow$ intent parsing $\rightarrow$ action execution $\rightarrow$ output sanitasi $\rightarrow$ grounding check.
    \item API Server mengembalikan respons ke WA Bot.
    \item WA Bot mengirimkan balasan ke pelanggan via WhatsApp.
    \item Dashboard Admin mengambil data dari API untuk ditampilkan dan dimodifikasi.
    \item Web-Gen membaca data toko dan produk dari API untuk menghasilkan website statis.
\end{enumerate}

\subsection{AI Pipeline (18 Langkah)}

\begin{enumerate}[nosep]
    \item \textbf{Normalisasi input} -- NFKC normalization, hapus karakter kontrol/zero-width, potong 4000 karakter.
    \item \textbf{Upsert pelanggan \& percakapan} -- Cari/buat Customer dan Conversation aktif.
    \item \textbf{Deduplikasi} -- Skip jika waMsgId sudah diproses.
    \item \textbf{Persistensi} -- Simpan Message dengan role = CUSTOMER.
    \item \textbf{Rate limit} -- Sliding window per pelanggan (8/30s + 60/jam).
    \item \textbf{PII scan} -- Deteksi nomor telepon, email, NIK, alamat.
    \item \textbf{3-tier injection defense} -- Lihat sub-bab terpisah.
    \item \textbf{Budget check} -- Verifikasi batas harian panggilan LLM (UsageCounter).
    \item \textbf{Load context} -- Store, Products, AiConfig.
    \item \textbf{Build prompt} -- System prompt + 10 pesan terakhir + current message dengan delimiters + canary token.
    \item \textbf{LLM completion} -- OpenRouter dengan circuit breaker, retry 2x, fallback model.
    \item \textbf{Parsing output} -- Ekstrak JSON, validasi Zod discriminated union (6 intent).
    \item \textbf{Handle intent} -- Eksekusi aksi: buat Order, eskalasi, dll.
    \item \textbf{Sanitasi reply} -- Hapus code fences, potong 1500 karakter.
    \item \textbf{Output scan} -- Deteksi kebocoran canary token, delimiter, system prompt, PII.
    \item \textbf{PII redaction} -- Ganti data pribadi dengan [PHONE], [EMAIL], dll.
    \item \textbf{Grounding check} -- Khusus inquiry/order: LLM judge verifikasi akurasi.
    \item \textbf{Record \& persist} -- Catat usage, simpan outbound message, update conversation.
\end{enumerate}

\subsection{Deteksi Intent}

Output LLM diparsing ke dalam enam kategori intent melalui Zod \textit{discriminated union}:

\begin{table}[H]
\centering
\caption{Enam Kategori Intent WANI}
\label{tab:intent}
\begin{tabularx}{\textwidth}{p{2.5cm}Xp{5cm}}
\toprule
\textbf{Intent} & \textbf{Deskripsi} & \textbf{Aksi Sistem} \\
\midrule
\texttt{order} & Pelanggan memesan produk & Buat entitas \texttt{Order} + \texttt{OrderItem} \\
\texttt{inquiry} & Menanyakan info produk & Balas dengan data dari katalog \\
\texttt{greeting} & Sapaan awal & Balas pesan sambutan konfigurabel \\
\texttt{complaint} & Keluhan & Balas empati, eskalasi jika diperlukan \\
\texttt{unknown} & Tidak terklasifikasi & Balas umum, minta klarifikasi \\
\texttt{escalate} & Butuh campur tangan manusia & Catat ke \texttt{ActivityLog}, notifikasi admin \\
\bottomrule
\end{tabularx}
\end{table}

\subsection{Pertahanan Injeksi 3-Tier}

\begin{table}[H]
\centering
\caption{Detail 3-Tier Injection Defense}
\label{tab:tier}
\begin{tabularx}{\textwidth}{cp{1.8cm}p{1.8cm}X}
\toprule
\textbf{Tier} & \textbf{Metode} & \textbf{Latency} & \textbf{Cakupan} \\
\midrule
T1 & Regex patterns & $\sim$0ms & 9 attack classes: delimiter escape, instruction override, prompt extraction, role hijack, authority claim, token injection, context overflow, crescendo marker, leet obfuscation. Dilengkapi homoglyph detector (Cyrillic, Greek, Letterlike Symbols) + leetspeak normalizer. \\
\addlinespace
T2 & ML Classifier & $\sim$500ms & OpenRouter fast model untuk kasus UNCERTAIN dari T1. \\
\addlinespace
T3 & Deep Judge & $\sim$1500ms & Analisis konteks percakapan penuh untuk kasus SUSPICIOUS dari T2. \\
\bottomrule
\end{tabularx}
\end{table}

% ============================================================
%  4. TEKNOLOGI YANG DIGUNAKAN
% ============================================================
\section{Teknologi yang Digunakan}

Seluruh teknologi yang digunakan dalam WANI adalah versi stabil terbaru dari masing-masing ekosistem, tidak ada yang diturunkan versinya:

\begin{table}[H]
\centering
\caption{Stack Teknologi WANI}
\label{tab:stack}
\begin{tabularx}{\textwidth}{p{2.8cm}p{3cm}X}
\toprule
\textbf{Kategori} & \textbf{Teknologi} & \textbf{Peran} \\
\midrule
Runtime & Bun 1.3+ & Runtime JavaScript/TypeScript cepat dengan \textit{built-in bundler}, test runner, package manager \\
Bahasa & TypeScript 5\textasciitilde6 & Static typing untuk keamanan dan ketertiban kode \\
\hline
Backend API & Express 5 & Web framework Node.js minimalis dengan middleware \\
ORM & Prisma 7 + \texttt{@prisma/adapter-pg} & ORM type-safe dengan database PostgreSQL \\
Database & PostgreSQL 16+ & Basis data relasional open-source \\
Validasi & Zod v4 & Validasi skema dengan inferensi tipe otomatis \\
Keamanan HTTP & Helmet 8 & Middleware keamanan Express \\
Logging & Winston 3 + Morgan & Logging multi-transport JSON \\
\hline
Frontend & React 19 + React Compiler & UI deklaratif dengan kompilasi otomatis \\
Build Tool & Vite 8 (Rolldown) & Build cepat dengan HMR \\
Styling & Tailwind CSS 4 & Utility-first CSS framework \\
Routing & React Router 8 & Client-side routing dengan loader/action \\
\hline
WhatsApp Bot & Baileys 6 & WhatsApp Web API multi-device \\
HTTP Client & Axios 1 & HTTP promise-based \\
Logging & Pino 9 & Logger berkinerja tinggi \\
\hline
Site Generator & Astro 6.4 & Static site generator dengan island architecture \\
ZIP & Archiver 7 & Pembuatan arsip ZIP \\
\hline
AI Engine & OpenRouter API & Agregator LLM (DeepSeek V4, Gemini 2.0) \\
Model Utama & DeepSeek V4 Flash & Chat completion dan intent parsing \\
Model Cadangan & Gemini 2.0 Flash & Fallback, classifier, judge, grounding \\
\bottomrule
\end{tabularx}
\end{table}

% ============================================================
%  5. FITUR DASHBOARD
% ============================================================
\section{Fitur Dashboard dan Antarmuka}

\subsection{Admin Dashboard}

Dashboard WANI terdiri dari 5 halaman utama yang saling terintegrasi, dibangun dengan React 19 dan Tailwind CSS 4:

\begin{enumerate}
    \item \textbf{Dashboard} -- Ringkasan statistik real-time: total produk, pesanan pending (dengan status), pelanggan aktif, indikator koneksi WhatsApp, dan daftar pesanan terbaru.
    \item \textbf{Products} -- Manajemen produk dengan dua mode tampilan (tabel dan kartu), pencarian, filter kategori, serta form tambah/edit produk dalam modal.
    \item \textbf{Orders} -- Daftar pesanan dengan filter status (Pending/Confirmed/Processing/Completed/Cancelled), detail pesanan lengkap dengan timeline status dan informasi pembayaran.
    \item \textbf{Customers} -- Panel ganda: daftar pelanggan di kiri dan percakapan \textit{inline chat} di kanan, memungkinkan admin untuk melihat riwayat dan membalas langsung.
    \item \textbf{Settings} -- Tiga tab konfigurasi: Store (profil bisnis), AI Agent (model LLM, prompt, temperature), WA Session (status koneksi dan QR code).
\end{enumerate}

\subsection{Generated Website}

Website yang dihasilkan oleh Web-Gen memiliki tiga halaman responsif dengan integrasi WhatsApp:

\begin{itemize}
    \item \textbf{Beranda} -- Hero section, tentang toko, dan produk unggulan.
    \item \textbf{Produk} -- Grid katalog produk dengan kartu produk dan tombol ``Pesan via WhatsApp'' yang otomatis mengisi nomor toko dan template pesanan.
    \item \textbf{Kontak} -- Informasi kontak, alamat, jam operasional, dan tautan langsung WhatsApp.
\end{itemize}

% ============================================================
%  6. METODOLOGI PENGEMBANGAN
% ============================================================
\section{Metodologi Pengembangan}

\subsection{Paradigma: Extreme Programming}

WANI dikembangkan menggunakan metodologi \textbf{Extreme Programming} (XP) dengan alasan:
\begin{itemize}
    \item Empat sub-proyek independen memungkinkan pengembangan paralel dan iteratif.
    \item \textit{Evolving requirements} yang khas dalam proyek inovasi.
    \item Fokus pada kualitas kode dan \textit{testing} yang ketat.
    \item Tim kecil (4 orang) ideal untuk \textit{pair programming} dan \textit{collective ownership}.
\end{itemize}

\subsection{Tahapan}

\begin{enumerate}
    \item \textbf{Planning} -- \textit{User stories} dari studi literatur dan observasi UMKM. Prioritas berdasarkan nilai bisnis. Iterasi 1--2 minggu.
    \item \textbf{Design} -- Prinsip YAGNI (\textit{You Ain't Gonna Need It}). Arsitektur modular dengan 13 entitas data. \textit{Entity Relationship Diagram} mencakup: MERCHANT, CUSTOMER, CATEGORY, PRODUCT, ORDER, ORDER\_ITEM, PAYMENT, CONVERSATION, MESSAGE, WA\_SESSION, AI\_AGENT, SETTING, ACTIVITY\_LOG.
    \item \textbf{Coding} -- \textit{Continuous integration}. Pekerjaan paralel per sub-proyek. API: arsitektur routes $\rightarrow$ controllers $\rightarrow$ models $\rightarrow$ Prisma delegate. Dashboard: SPA dengan React Router v8. Bot: Baileys 6 dengan auth persisten PostgreSQL. Web-Gen: Astro 6.4 dengan template Tailwind.
    \item \textbf{Testing} -- TDD dengan bun:test. Cakupan: unit test (Zod, utilitas, regex), integration test (AI pipeline, circuit breaker, rate limiter), security test (3-tier firewall, PII, grounding), golden reply test.
    \item \textbf{Release} -- Siklus pendek, rilis fitur stabil, dokumentasi changelog.
\end{enumerate}

% ============================================================
%  7. PROFIL TIM
% ============================================================
\section{Profil Tim Pengembang}

Tim WANI terdiri dari empat mahasiswa Program Studi Teknik Informatika Universitas Darussalam Gontor dengan pembagian peran sebagai berikut:

\begin{table}[H]
\centering
\caption{Profil Tim Pengembang}
\label{tab:tim}
\begin{tabularx}{\textwidth}{p{2.5cm}p{4cm}X}
\toprule
\textbf{Peran} & \textbf{Nama} & \textbf{Kompetensi} \\
\midrule
\textbf{Ketua Tim} / & Farrel Ghozy & Full-Stack Web Developer, \\
\textit{Backend \& AI} & Affifudin & TypeScript, Node.js, Express, \\
 & & React, Prisma, PostgreSQL, \\
 & & Docker, Linux Server \\
\hline
Anggota 1 / & [Nama Anggota 1] & [Kompetensi] \\
\textit{Frontend} & & \\
\hline
Anggota 2 / & [Nama Anggota 2] & [Kompetensi] \\
\textit{WhatsApp Bot} & & \\
\hline
Anggota 3 / & [Nama Anggota 3] & [Kompetensi] \\
\textit{Web-Gen} & & \\
\hline
\end{tabularx}
\end{table}

\noindent
\textbf{Profil Ketua Tim:}
\begin{itemize}
    \item Mahasiswa Teknik Informatika, Universitas Darussalam Gontor (Semester 5, IPK 3.79/4.00).
    \item Berpengalaman dalam pengembangan \textit{full-stack} web, \textit{system administration} (Ubuntu, Docker, Nginx, Proxmox, K3s), dan \textit{machine learning} (TensorFlow, PyTorch).
    \item Prestasi: Juara 2 Coding Day (2026), The Best Thematic Hackathon SRE ITB (2026).
    \item Ketua Himpunan Mahasiswa Teknik Informatika (2025--sekarang).
\end{itemize}

% ============================================================
%  8. RENCANA PENGEMBANGAN
% ============================================================
\section{Rencana Pengembangan}

\subsection{Capaian Saat Ini}

\begin{itemize}
    \item \textbf{API Server}: 100\% -- 6 endpoint REST, Zod validation, Winston logging, Prisma ORM, graceful shutdown.
    \item \textbf{Dashboard}: 100\% UI -- 5 halaman, routing, komponen UI, hooks dengan MOCK toggle.
    \item \textbf{WA Bot}: 100\% -- Baileys 6, QR relay, auto-reconnect, auth persisten via Prisma.
    \item \textbf{AI Pipeline}: 100\% -- 18 langkah, 3-tier fire wall, 6 intent detection, circuit breaker, grounding.
    \item \textbf{Guardrails}: 100\% -- PII scanner, rate limiter, budget tracker, leetspeak/homoglyph detector.
    \item \textbf{Web-Gen}: 100\% -- template Astro, generator, ZIP archiver.
    \item \textbf{Testing}: 4 file test (guardrail, firewall, intent, golden reply).
    \item \textbf{Bot $\rightarrow$ API integrasi}: End-to-end connected.
\end{itemize}

\subsection{Rencana ke Depan}

\begin{table}[H]
\centering
\caption{Roadmap Pengembangan WANI}
\label{tab:roadmap}
\begin{tabularx}{\textwidth}{p{2cm}p{4cm}X}
\toprule
\textbf{Fase} & \textbf{Fitur} & \textbf{Keterangan} \\
\midrule
Fase 1 & API Endpoints Produk, & REST CRUD untuk manajemen \\
(Juli 2026) & Orders, Customers & data lengkap \\
\hline
Fase 2 & API Endpoints Web-Gen, & Preview ZIP, publish, \\
(Agustus 2026) & Deployment & static serving \\
\hline
Fase 3 & Multi-store support, & Tenant management, \\
(September 2026) & Role-based access & autentikasi multi-level \\
\hline
Fase 4 & RAG / Vector database, & Embeddings untuk knowledge \\
(Oktober 2026) & Evaluasi harness & retrieval, benchmark AI \\
\hline
\end{tabularx}
\end{table}

% ============================================================
%  9. PANDUAN INSTALASI
% ============================================================
\section{Panduan Instalasi}

\subsection{Prasyarat Sistem}

\begin{itemize}
    \item \textbf{Bun 1.3+} -- Runtime JavaScript (\url{https://bun.sh})
    \item \textbf{PostgreSQL 16+} -- Basis data relasional
    \item \textbf{Git} -- Version control
\end{itemize}

\subsection{Langkah Instalasi}

\begin{enumerate}
    \item \textbf{Clone repositori}
\begin{lstlisting}[style=bash]
git clone https://github.com/FarrelGhozy/WANI.git
cd WANI
\end{lstlisting}

    \item \textbf{Install dependencies}
\begin{lstlisting}[style=bash]
cd api && bun install
cd ../dashboard && bun install
cd ../web-gen && bun install
cd ../wa-bot && bun install
\end{lstlisting}

    \item \textbf{Buat database PostgreSQL}
\begin{lstlisting}[style=bash]
createdb -U postgres wani_api
createdb -U postgres wa_bot
\end{lstlisting}

    \item \textbf{Konfigurasi environment variable} \\
    Salin \texttt{.env.example} jadi \texttt{.env} di \texttt{api/}, \texttt{wa-bot/}, \texttt{web-gen/}. \\
    Variabel penting:
    \begin{itemize}
        \item \texttt{DATABASE\_*} -- koneksi PostgreSQL
        \item \texttt{API\_TOKEN} -- shared secret autentikasi bot $\leftrightarrow$ API
        \item \texttt{OPENROUTER\_API\_KEY} -- kunci API untuk LLM
    \end{itemize}

    \item \textbf{Jalankan migrasi database}
\begin{lstlisting}[style=bash]
cd api && bunx --bun prisma migrate dev
cd ../wa-bot && bunx --bun prisma migrate dev
\end{lstlisting}

    \item \textbf{Jalankan service} (API harus sebelum Bot)
\begin{lstlisting}[style=bash]
# Terminal 1 -- API Server (port 3001)
cd api && bun run src/index.ts

# Terminal 2 -- WhatsApp Bot
cd wa-bot && bun run src/index.ts

# Terminal 3 -- Dashboard (port 5173)
cd dashboard && bun run dev

# Terminal 4 -- Web-Gen (testing)
cd web-gen && bun run src/generator.ts --slug test \
    --template default --data ./test-data.json
\end{lstlisting}
\end{enumerate}

\subsection{Verifikasi}

\begin{itemize}
    \item Dashboard: \url{http://localhost:5173}
    \item API Server: \url{http://localhost:3001}
    \item Status WhatsApp: \texttt{GET /api/qr/status}
\end{itemize}

% ============================================================
%  10. KESIMPULAN
% ============================================================
\section{Kesimpulan}

WANI (\textit{WhatsApp AI Native Integration} / \textit{WhatsApp Niaga}) menawarkan solusi omnichannel yang komprehensif, aman, dan terjangkau bagi pelaku UMKM Indonesia. Dengan mengintegrasikan WhatsApp Bot bertenaga AI, REST API server, admin dashboard, dan website generator dalam satu platform open-source, WANI menjawab tiga masalah utama UMKM:

\begin{enumerate}
    \item \textbf{Ketidakefisienan operasional} -- melalui otomatisasi layanan pelanggan 24/7 dengan AI chatbot yang mampu mendeteksi dan merespons 6 kategori intensi pelanggan.
    \item \textbf{Kerentanan keamanan AI} -- melalui sistem pertahanan injeksi 3-tier yang menjadi standar keamanan tertinggi untuk chatbot AI di lingkungan UMKM.
    \item \textbf{Ketiadaan web presence} -- melalui website generator otomatis yang menghasilkan situs web statis lengkap dengan integrasi WhatsApp.
\end{enumerate}

\noindent
Dibangun di atas teknologi modern dan stabil (Bun, Express 5, Prisma 7, React 19, Astro 6.4) dengan metodologi Extreme Programming, WANI telah siap untuk diimplementasikan dan dikembangkan lebih lanjut. Seluruh kode sumber tersedia secara \textit{open source} di \href{https://github.com/FarrelGhozy/WANI}{GitHub}, memungkinkan kontribusi dan adopsi yang luas oleh komunitas.

% ============================================================
%  DAFTAR PUSTAKA
% ============================================================
\newpage
\begin{thebibliography}{99}
\addcontentsline{toc}{section}{Daftar Pustaka}

\bibitem{baileys2024}
WhiskeySockets. (2024). \textit{Baileys: WhatsApp Web API Library for Node.js}. GitHub Repository. \\
\url{https://github.com/WhiskeySockets/Baileys}

\bibitem{express2024}
OpenJS Foundation. (2024). \textit{Express.js -- Node.js Web Application Framework}. \\
\url{https://expressjs.com}

\bibitem{prisma2024}
Prisma Data, Inc. (2024). \textit{Prisma ORM Documentation}. \\
\url{https://www.prisma.io/docs}

\bibitem{react2024}
Meta Platforms, Inc. (2024). \textit{React Documentation}. \\
\url{https://react.dev}

\bibitem{vite2024}
VoidZero Inc. (2024). \textit{Vite Documentation}. \\
\url{https://vite.dev}

\bibitem{astro2024}
Astro. (2024). \textit{Astro Documentation}. \\
\url{https://docs.astro.build}

\bibitem{tailwind2024}
Tailwind Labs. (2024). \textit{Tailwind CSS Documentation}. \\
\url{https://tailwindcss.com/docs}

\bibitem{zod2024}
Colin McDonnell. (2024). \textit{Zod -- TypeScript-first Schema Validation}. \\
\url{https://zod.dev}

\bibitem{openrouter2024}
OpenRouter. (2024). \textit{OpenRouter Documentation}. \\
\url{https://openrouter.ai/docs}

\bibitem{bun2024}
Oven-Sh. (2024). \textit{Bun -- JavaScript Runtime Documentation}. \\
\url{https://bun.sh/docs}

\bibitem{typescript2024}
Microsoft. (2024). \textit{TypeScript Documentation}. \\
\url{https://www.typescriptlang.org/docs}

\bibitem{helmet2024}
Evan Hahn. (2024). \textit{Helmet -- Secure Express Apps}. \\
\url{https://helmetjs.github.io}

\bibitem{winston2024}
Charlie Robbins. (2024). \textit{Winston -- A Logger for Node.js}. \\
\url{https://github.com/winstonjs/winston}

\bibitem{pino2024}
Pino Contributors. (2024). \textit{Pino -- Very Low Overhead Node.js Logger}. \\
\url{https://github.com/pinojs/pino}

\bibitem{axios2024}
Axios Contributors. (2024). \textit{Axios -- Promise-based HTTP Client}. \\
\url{https://axios-http.com}

\bibitem{reactrouter2024}
Shopify Inc. (2024). \textit{React Router Documentation}. \\
\url{https://reactrouter.com}

\bibitem{deepseek2024}
DeepSeek. (2024). \textit{DeepSeek V4 Flash Model}. OpenRouter. \\
\url{https://openrouter.ai/models/opencode/deepseek-v4-flash-free}

\bibitem{gemini2024}
Google DeepMind. (2024). \textit{Gemini 2.0 Flash}. OpenRouter. \\
\url{https://openrouter.ai/models/google/gemini-2.0-flash-exp:free}

\bibitem{archiver2024}
Chris Talkington. (2024). \textit{Archiver -- Streaming Archive Library for Node.js}. \\
\url{https://github.com/archiverjs/node-archiver}

\bibitem{postgres2024}
PostgreSQL Global Development Group. (2024). \textit{PostgreSQL Documentation}. \\
\url{https://www.postgresql.org/docs}

\bibitem{owasp2025}
OWASP. (2025). \textit{OWASP Top 10 for LLM Applications v2.0}. \\
\url{https://genai.owasp.org/llm-top-10/}

\bibitem{kemkop2023}
Kementerian Koperasi dan UKM. (2023). \textit{Kontribusi UMKM terhadap PDB Indonesia}. \\
\url{https://djpb.kemenkeu.go.id/kppn/lubuksikaping/id/data-publikasi/artikel/3134-kontribusi-umkm-dalam-perekonomianindonesia.html}

\bibitem{kemkop2022survei}
Kementerian Koperasi dan UKM. (2022). Survei Adopsi Teknologi Digital UMKM. Dalam: Medcom.id. \\
\url{https://www.medcom.id/ekonomi/ekonomi-digital/GbmM1zyb-adopsi-digital-perlu-dilakukan-pelaku-umkm}

\bibitem{wearesocial2024}
We Are Social. (2024). \textit{Digital 2024: Indonesia}. \\
\url{https://wearesocial.com/id/blog/2024/01/digital-2024/}

\end{thebibliography}

% ============================================================
%  LAMPIRAN
% ============================================================
\newpage
\section*{Lampiran}
\addcontentsline{toc}{section}{Lampiran}

\subsection*{Lampiran 1: Tautan Repositori Kode}

\begin{itemize}
    \item \textbf{Repositori GitHub (Open Source):} \url{https://github.com/FarrelGhozy/WANI}
\end{itemize}

\noindent
Repositori ini mencakup seluruh kode sumber WANI: \texttt{api/}, \texttt{dashboard/}, \texttt{wa-bot/}, \texttt{web-gen/}, beserta dokumentasi dan konfigurasi.

\subsection*{Lampiran 2: Tautan Live Demo / Production Deployment}

\begin{itemize}
    \item \textbf{Live Demo}: \texttt{[URL Live Demo -- belum tersedia]}
\end{itemize}

\subsection*{Lampiran 3: Tautan Unduhan Berkas Aplikasi}

Tidak berlaku. WANI merupakan platform berbasis \textit{website} yang berjalan di atas server, bukan aplikasi \textit{mobile}. Tidak ada berkas \texttt{.apk}.

\subsection*{Lampiran 4: Tautan Prototipe Interaktif (Figma)}

Tidak ada. Desain antarmuka WANI dikembangkan langsung dalam bentuk kode (\textit{in-code design}) menggunakan React 19 dan Tailwind CSS 4, sejalan dengan metodologi Extreme Programming yang mengutamakan iterasi cepat.

\end{document}
